%% =============================================================================
%% sm_c1_pseudo.tex — SM-C, Subsection C.1: Pseudo-Likelihood and Separability
%% Body content only (\subsection level); parent structure in sm_c.tex
%% Primary source: dev/log/06_research-notes-survey-design-pseudo-posterior.qmd
%%   lines 570--660 (normalization, separability, weighted scores)
%% =============================================================================

\subsection{Pseudo-likelihood and separability}
\label{smc:pseudo-lik}

We summarize the pseudo-likelihood structure
from~\cref{sec:survey}.
By~\cref{prop:weighted-sep}, the decomposition
$\ell^{(w)}
 = \ell_{\mathrm{ext}}^{(w)}
 \allowbreak+ \ell_{\mathrm{int}}^{(w)}$
holds with weights
$\tilde{w}_i = w_i N / \sum_j w_j$.
We record three clarifying remarks.

\begin{remark}[Weight normalization]\label{smc:rem-normalization}
  Two normalization conventions appear in the survey literature.
  The \emph{sum-to-$N$} convention, $\tilde{w}_i = w_i N / \sum_j w_j$
  so that $\sum_i \tilde{w}_i = N$, is adopted throughout and is standard
  in the Bayesian pseudo-posterior literature
  \citep{SavitskyToth2016,WilliamsSavitsky2021}.
  The alternative \emph{sum-to-$\hat{M}$} convention retains unnormalized
  weights $\tilde{w}_i = w_i$, which for the NSECE gives
  $\sum_i w_i \approx 119{,}593$.  Under this convention the
  pseudo-likelihood resembles roughly $119{,}593$ observations rather
  than $N = 6{,}809$, leading to extreme overconcentration of the
  pseudo-posterior and credible intervals that are far too narrow.
  The sum-to-$N$ normalization ensures that the pseudo-likelihood is
  comparable in scale to an ordinary likelihood based on $N$ observations.
\end{remark}

\begin{remark}[Weighted score functions]\label{smc:rem-scores}
  The observation-level score functions for the extensive and intensive
  margins, along with their Fisher information and sign verification,
  are developed in~\cref{smb:scores} (Supplementary Material~B).
  The weighted versions used here simply multiply each per-observation
  score by $\tilde{w}_i$; the separability of
  \cref{prop:weighted-sep} guarantees that the weighted
  extensive-margin scores depend only on $(\balpha, \bdelta_{1,\cdot})$
  and the weighted intensive-margin scores depend only on
  $(\bbeta, \bdelta_{2,\cdot}, \kappa)$.
\end{remark}

\begin{remark}[Significance of separability under weighting]
\label{smc:rem-separability-significance}
  The separability result in~\cref{prop:weighted-sep} is not
  guaranteed for all survey adjustments.  Approaches that modify the
  likelihood structure---such as multilevel weight smoothing
  \citep{RabeHeskethSkrondal2006} or conditional likelihood
  methods---can break the additive structure.  Observation-level
  exponentiation $[f_i]^{\tilde{w}_i}$ preserves it precisely because
  $\tilde{w}_i$ is a scalar multiplier in the log scale and does
  not couple the two parameter blocks.  This separability is exploited
  in the sandwich estimator (\cref{sec:survey}), where the
  block-diagonality of $\mathbf{H}_{\mathrm{obs}}$ follows directly
  from the additive decomposition.
\end{remark}
